\source{Лекции, с. 15–21; Карпов, §1.5.2, с. 25–27}
\begin{definition}[Лес, дерево, лист]
Ациклический неориентированный граф (без простых циклов) называется лесом.
Дерево — связный ациклический граф. Лист — вершина степени 1.
Остовное дерево — дерево, являющееся остовным подграфом исходного графа.
\end{definition}
\begin{lemma}[Листья]
В дереве с не менее чем двумя вершинами есть хотя бы два листа.
Удаление листа с его ребром оставляет дерево.
\end{lemma}
\begin{proof}
Возьмём простой путь максимальной длины. У каждого его конца нет соседа
вне пути (путь удлинился бы) и нет другого соседа на пути, кроме соседней
вершины (возник бы цикл). Оба конца — листья.
Путь между оставшимися вершинами не проходит через лист внутренней
вершиной, поэтому после удаления связность сохраняется; циклов не возникает.
\end{proof}
Для одноточечного дерева листьев нет; утверждения о листьях требуют \(n\ge2\).
\begin{theorem}[Любые два из трёх свойств]
Для непустого графа любые два из следующих свойств влекут третье:
\begin{enumerate}[nosep]
\item \(m=n-1\);
\item граф связен;
\item граф ацикличен.
\end{enumerate}
Следовательно, любая пара характеризует дерево.
\end{theorem}
\begin{proof}
\(2,3\Rightarrow1\). Индукция по \(n\): при \(n=1\) рёбер нет.
При \(n>1\) удаляем лист. Для оставшегося дерева \(m'=n'-1\),
поэтому \(m=m'+1=(n-1)-1+1=n-1\).

\(1,3\Rightarrow2\). Каждая из \(c\) компонент — дерево.
Если размеры компонент \(n_1,\ldots,n_c\), по уже доказанному
\(m=\sum_i(n_i-1)=n-c\). Сравнивая с \(m=n-1\), получаем \(c=1\).

\(1,2\Rightarrow3\). Снова индукция по \(n\). При \(n=1\) очевидно.
При \(n>1\) степени положительны, а их сумма \(2n-2<2n\),
поэтому есть вершина степени 1. Удаление этого листа сохраняет связность:
простой путь между оставшимися вершинами не мог использовать лист внутри.
Остаётся связный граф на \(n-1\) вершинах с \(n-2\) рёбрами; по индукции
он ацикличен. Добавление листа цикла не создаёт.
\end{proof}
\begin{cor}[Число рёбер леса]\label{lem:forest}
Для леса \(m=n-c\); для связного графа \(m\ge n-1\).
\end{cor}
\begin{proof}
Первое равенство уже получено суммированием по компонентам.
Если связный граф имеет цикл, удаляем из него ребро, сохраняя связность.
Повторяем до дерева. В нём \(n-1\) рёбер, поэтому исходно их не меньше.
\end{proof}
\begin{theorem}[Единственность пути]
Граф — дерево тогда и только тогда, когда между любыми двумя различными
вершинами существует ровно один простой путь.
\end{theorem}
\begin{proof}
В дереве связность даёт путь. У двух разных простых путей возьмём
первое расхождение и следующую встречу: участки между ними образуют цикл.
Обратно, условие даёт связность, а простой цикл дал бы два пути
между любыми двумя своими различными вершинами.
\end{proof}
\begin{theorem}[Минимальная связность и максимальная ацикличность]
Для непустого графа эквивалентны: быть деревом; быть связным и терять
связность при удалении любого ребра; быть ацикличным и получать цикл
при добавлении любого отсутствующего ребра.
\end{theorem}
\begin{proof}
В дереве каждое ребро — мост, а добавление ребра замыкает единственный
путь между его концами в цикл. Обратно, минимально связный граф не может
иметь цикл, поскольку его ребро удалялось бы без потери связности.
Максимально ацикличный граф не может иметь две компоненты: ребро между
ними цикла не создаёт. В обоих случаях получено дерево.
\end{proof}
\begin{cor}[Вершины и остов]
В дереве каждая вершина степени больше 1 является точкой сочленения.
Каждый связный граф имеет остовное дерево.
\end{cor}
\begin{proof}
После удаления нелистовой вершины её разные соседи не могут остаться
связанными: такой обход дал бы цикл в дереве. Существование остова
доказывается последовательным удалением рёбер циклов, как выше.
\end{proof}
\textbf{Уточнение записи с. 19.} Обратное утверждение о точках сочленения
не используется: треугольник с листом у каждой вершины имеет все нелистовые
вершины точками сочленения, но деревом не является.
